\section{SNAC Audio Processing Pipeline Implementation}
\label{app:snac_pipeline}
To enable large-scale audio tokenization for multimodal training, we used an optimized SNAC processing pipeline that addresses memory constraints and variable-length audio challenges through adaptive bucketing and batch processing.

\subsection{Technical Implementation}
Our pipeline groups audio files into duration-based buckets ($\leq$15s, 15-45s, $>$45s) with corresponding target lengths (10s, 30s, 60s). Within each bucket, audio sequences undergo truncation or zero-padding to achieve uniform dimensions, enabling efficient GPU batch processing while preventing out-of-memory errors.

The output maintains temporal structure through 7-token groupings corresponding to SNAC's internal resolution:

\begin{verbatim}
{
  "id": "videoID_chunkIDX",
  "text": "corresponding transcript text",
  "snac_tokens": [
    [128266, 128267, 169362, 171758, 210154, 250239, 290518],
    [128268, 128269, 169363, 171759, 210155, 250240, 290519],
    [...]
  ],
  "group_duration_sec": 0.018375
}
\end{verbatim}

This structured format preserves precise temporal alignment (0.018375s per group from hop\_length=441 at 24kHz) essential for downstream speech synthesis tasks.

\subsection{Performance Results}
The implementation achieved:
\begin{itemize}
\item Memory-efficient processing through adaptive bucketing
\item Scalable batch processing leveraging GPU parallelization  
\item Curently used for processing of large-scale datasets (tested on ontocord/VALID)
\item Maintained high-fidelity token representations with temporal precision
\end{itemize}
